\documentclass[12pt]{article}
\usepackage{amsmath}
\usepackage{amssymb}
\usepackage{geometry}
\geometry{a4paper, margin=1in}

\begin{document}

\begin{center}
    \textbf{\Large Kontes Terbuka Olimpiade Matematika - Simulasi OSN-K SMP 2025}
\end{center}

\vspace{0.5cm}

\begin{enumerate}

    % answer  : D
    \item Tentukan banyaknya bilangan asli $1 \le n \le 50$ sehingga 
    $\frac{1\times2\times3\times\cdots\times n}{1+2+3+\cdots+n}$ 
    merupakan bilangan bulat.
    \begin{enumerate}
        \item 33
        \item 34
        \item 35
        \item 36
    \end{enumerate}

    % answer  : D
    \item Diketahui bahwa bilangan riil $x, y$ memenuhi persamaan $x^{2}+y^{2}=10x+8y+5$. Tentukan bilangan asli terkecil yang lebih besar dari $5x+2y$ untuk seluruh kemungkinan $(x, y)$.
    \begin{enumerate}
        \item 67
        \item 68
        \item 69
        \item 70
    \end{enumerate}

    % answer  : C
    \item Diberikan segitiga sama kaki $ABC$ dengan $AB=AC$ dan $\angle A=20^{\circ}$. Titik $P$ terletak di sisi $AC$ sehingga $AP=BC$. Tentukan besar $\angle PBC$.
    \begin{enumerate}
        \item 60
        \item 65
        \item 70
        \item 75
    \end{enumerate}

    % answer  : B
    \item Jika $x$ ialah bilangan real sehingga $x+\frac{1}{x}=\sqrt{7}$, tentukan nilai dari $\lfloor x^{9}-23x^{5}+x \rfloor$.
    \begin{enumerate}
        \item -1
        \item 0
        \item 1
        \item 2
    \end{enumerate}

    % answer  : B
    \item Tentukan banyaknya bilangan asli $n$ unik sehingga terdapat bilangan bulat positif $x \in \{1,2,3,\dots,2000\}$ yang memenuhi $\lceil\frac{x^{2}+2x}{2025}\rceil=n$.
    \begin{enumerate}
        \item 495
        \item 1495
        \item 1395
        \item 395
    \end{enumerate}

    % answer  : A
    \item Diketahui segitiga $ABC$ dan titik $D$ adalah titik pada lingkaran luar $ABC$ sehingga $BD$ adalah garis yang membagi sudut $\angle ABC$ menjadi dua sudut yang sama besar. Titik $E$ dan $F$ berturut-turut adalah perpotongan $AB$ dan $BC$ dengan garis yang melewati $D$ dan menyinggung lingkaran luar $ABC$. Jika $AC=7$, $BC=8$, dan $AB=6$ maka $DE \cdot DF=\frac{a}{b}$ dengan $\text{FPB}(a,b)=1$. Tentukan nilai $a+b$.
    \begin{enumerate}
        \item 67
        \item 84
        \item 57
        \item 71
    \end{enumerate}

    % answer  : D
    \item Tentukan hasil kali dari seluruh bilangan bulat $k$ sehingga persamaan kuadrat $x^{2}-2kx+7k+2=0$ memiliki akar bilangan bulat.
    \begin{enumerate}
        \item 162
        \item 324
        \item 1444
        \item 3564
    \end{enumerate}

    % answer  : A
    \item Diberikan 2 kelompok data nilai siswa kelas A dan B. Jika rata-rata nilai kelas A dan B berturut-turut adalah 68 dan 83, tentukan banyak total siswa minimal di kelas A dan B sehingga rata-rata nilai gabungan kedua kelas adalah 79.
    \begin{enumerate}
        \item 15
        \item 16
        \item 30
        \item 32
    \end{enumerate}

    % answer  : B
    \item Sebuah keranjang berisi 6 telur ayam dan 12 telur bebek. Budi akan mengambil 2 telur ayam untuk makan malam, tetapi Budi tidak dapat melihat isi keranjangnya. Diketahui bahwa Budi hanya akan mengambil 4 telur saja dari keranjang tanpa pengembalian. Jika peluang Budi berhasil mendapatkan 2 telur ayam sebelum pengambilan keempat adalah $\frac{a}{b}$ di mana $a, b$ relatif prima, tentukan nilai dari $a+b$.
    \begin{enumerate}
        \item 126
        \item 127
        \item 128
        \item 256
    \end{enumerate}

    % answer  : D
    \item Andi memilih sebuah titik yang terletak pada parabola $y=x^{2}-6x+11$ dan berada di daerah $x \ge y-11$. Jika luas maksimum dari segitiga yang terbentuk dari 2 titik perpotongan garis $y=x+11$ dengan parabola dan titik yang dipilih Andi dapat dinyatakan sebagai $\frac{a}{b}$ dengan $a, b$ relatif prima, tentukan nilai $a+9b$.
    \begin{enumerate}
        \item 115
        \item 215
        \item 315
        \item 415
    \end{enumerate}

    % answer  : D
    \item Tentukan jumlah semua faktor prima dari $144^{3}-125^{2}$.
    \begin{enumerate}
        \item 56
        \item 169
        \item 273
        \item 362
    \end{enumerate}

    % answer  : B
    \item Berapakah banyak cara membagi bilangan $\{1,2,\dots,18\}$ menjadi 9 pasang bilangan sehingga untuk setiap pasangan, bilangan terbesar minimal dua kali lebih besar dari bilangan terkecil?
    \begin{enumerate}
        \item 1440
        \item 2880
        \item 4320
        \item 5760
    \end{enumerate}

    % answer  : B
    \item Tentukan banyak bilangan asli $n \le 2025$ sehingga persamaan $(n+1)x+(n+5)y=n^{2}+3n+5$ memiliki setidaknya satu solusi bilangan bulat.
    \begin{enumerate}
        \item 506
        \item 1012
        \item 1518
        \item 2024
    \end{enumerate}

    % answer  : C
    \item Diberikan barisan dengan $a_{1}=3$, $a_{2}=7$, dan $a_{n}=a_{n-1}+2a_{n-2}$. Tentukan banyak bilangan kubik yang habis membagi $a_{2025}+a_{2024}$.
    \begin{enumerate}
        \item 673
        \item 674
        \item 675
        \item 676
    \end{enumerate}

    % answer  : B
    \item Diberikan persegi $ABCD$ dengan titik $E$ dan $F$ terletak pada diagonal $AC$ sehingga $AE=EF=FC$. Titik $G$ dan $H$ terletak pada diagonal $BD$ sehingga $BG=GH=HD$. Jika panjang sisi persegi $ABCD$ adalah 9, tentukan luas dari segiempat $EHFG$.
    \begin{enumerate}
        \item 8
        \item 9
        \item 6
        \item 18
    \end{enumerate}

    % answer  : B
    \item Tentukan jumlah semua bilangan prima yang kurang dari 20 dan membagi $131313\dots13$ (dimana ada bilangan 13 muncul sebanyak 16 kali).
    \begin{enumerate}
        \item 24
        \item 30
        \item 32
        \item 36
    \end{enumerate}

    % answer  : C
    \item Misalkan $N=3^{2025}+4^{2025}$. Tentukan dua digit terakhir dari $N$.
    \begin{enumerate}
        \item 27
        \item 47
        \item 67
        \item 87
    \end{enumerate}

    % answer  : A
    \item Diberikan sepuluh kotak yang masing-masing berisi kelereng sebanyak 1, 2, 3, 4, 5, 6, 7, 8, 9, dan 10. Alice, Bob, dan Charlie masing-masing harus memilih tepat satu kotak dari sepuluh kotak tersebut dengan syarat tidak ada orang yang memilih kotak yang sama. Peluang Alice mendapatkan kelereng paling banyak dapat dinyatakan dalam bentuk $\frac{a}{b}$ di mana $a, b$ bilangan asli dan $\text{FPB}(a,b)=1$. Tentukan nilai dari $10a+b$.
    \begin{enumerate}
        \item 13
        \item 18
        \item 23
        \item 27
    \end{enumerate}

    % answer  : D
    \item Misalkan $P$ merupakan hasil perkalian akar-akar non-riil dari $x^{4}-4x^{3}+6x^{2}-4x=2025$. Tentukan $\lfloor P \rfloor$.
    \begin{enumerate}
        \item 46
        \item 44
        \item 44
        \item 46
    \end{enumerate}

    % answer  : A
    \item Perhatikan gambar berikut: (gunakan link berikut untuk mengakses gambar: https://acesse.one/DpnBb). Pada bangun A, tinggi kerucut 2 kali lipat tinggi tabung. Pada bangun B, bola menyentuh alas dan selimut tabung. Diketahui bahwa Bangun A dan Bangun B memiliki tinggi dan jari-jari yang sama. Jika volume gabungan pada bangun A sama dengan volume tabung yang diluar bola pada bangun B, tentukan perbandingan jari-jari dan tinggi bangun tersebut.
    \begin{enumerate}
        \item $\frac{1}{3}$
        \item $\frac{1}{4}$
        \item 3
        \item 4
    \end{enumerate}

    % answer  : B
    \item Tentukan banyak tripel bilangan bulat $(a, b, c)$ yang memenuhi sistem persamaan berikut: 
    $ab+c=26$, $bc+a=34$, $a+b+c=15$.
    \begin{enumerate}
        \item 0
        \item 1
        \item 2
        \item 3
    \end{enumerate}

    % answer  : D
    \item Untuk merayakan hari kemerdekaan NKRI, panitia menyediakan empat buah pot dengan motif yang berbeda untuk mengadakan lomba tabur benih bunga. Setiap peserta diberikan 12 benih identik untuk ditabur ke pot secara acak. Peserta dianggap lolos ke babak final jika berhasil mengisi paling sedikit 2 benih pada tepat 2 buah pot. Jika peluang Tika lolos ke babak final dapat dinyatakan sebagai pecahan paling sederhana $\frac{p}{q}$ tentukan nilai $p+q$.
    \begin{enumerate}
        \item 487
        \item 568
        \item 612
        \item 647
    \end{enumerate}

    % answer  : C
    \item Diberikan dua jenis ubin seperti pada gambar di bawah ini yang akan digunakan untuk menutup lantai. (gunakan link berikut untuk mengakses gambar: https://link.dev/N6LLc). Tentukan banyak cara menutup lantai berukuran $2\times7$ dengan syarat tidak ada ubin yang tumpang tindih dan ubin boleh dirotasi.
    \begin{enumerate}
        \item 645
        \item 650
        \item 655
        \item 660
    \end{enumerate}

    % answer  : C
    \item Misalkan $A$ menyatakan himpunan $\{1,2,3,4,5,6,7,8,9\}$. Berapa banyak fungsi $f:A\rightarrow A$ sehingga $f(f(f(x)))=x$?
    \begin{enumerate}
        \item 1800
        \item 2109
        \item 2409
        \item 2682
    \end{enumerate}

    % answer  : -
    \item Diketahui panjang 9 jembatan yang akan dibangun adalah bilangan bulat. Rata-rata dari 5 jembatan terpendek adalah 8, sementara rata-rata dari 6 jembatan terpanjang adalah 10. Jika rata-rata dari panjang 9 jembatan tersebut adalah 9, tentukan jumlah dari seluruh panjang yang mungkin dari jembatan terpanjang.
    \begin{enumerate}
        \item 77
        \item 87
        \item 88
        \item 98
    \end{enumerate}

\end{enumerate}

\end{document}